%%%%%%%%%%%%%%%%%%%%%%%%%%%%%%%%%%%%%%%%%%%%%%%%%%%%
%    Harvard Data Science Review Latex Template    %
%%%%%%%%%%%%%%%%%%%%%%%%%%%%%%%%%%%%%%%%%%%%%%%%%%%%

\documentclass[]{hdsr}

%Graphics should all go in the figs/ directory
\graphicspath{{figs/}}

\usepackage{xcolor}
\usepackage{hyperref}
\usepackage{cleveref}

\begin{document}

% Larger bottom margin for the first page
\newgeometry{bottom=1.5in}

% Editorial staff will replace the following values:
% 1. Volume number
% 2. Issue number
% 3. Article DOI
% e.g. for Volume 2, Issue 3, DOI 12.345:
% \volumeheader{2}{3}{12.345}
\volumeheader{0}{0}{00.000}



\begin{center}

  \title{Article Proposal: Portraying the Many Facets of Impact with Influence Flowers}
  \maketitle

  % Start page numbering on second page. Must appear *after* \maketitle
  \thispagestyle{empty}
  
  \vspace*{.2in}

  % Authors and Affiliations
  \begin{tabular}{cc}
    Minjeong Shin\upstairs{1}, Lexing Xie\upstairs{1}
   \\[0.25ex]
   {\small \upstairs{1} School of Computing, The Australian National University, Canberra, Australia} \\
  \end{tabular}
  
  % Replace with corresponding author email address
  % \emails{
  %   \upstairs{*}correspondingauthor@example.edu 
  %   }
  % \vspace*{0.4in}

\begin{figure}[H]
    \centering
    \includegraphics[width=\textwidth]{figs/hdsr_flower.pdf} \\
    \caption{\textbf{An influence flower visualization for HDSR (2019 to 2024)}, an {\href{https://influencemap.cmlab.dev/submit/?id=BAAAAAEAsCg5}{interactive version}} is available online. \textcolor{blue}{Blue arcs} denote incoming influence from the authors to papers in HDSR, with their thickness proportional to the number of references made. \textcolor{red}{Red arcs} denote outgoing influence from to papers in HDSR to the authors, with their thickness proportional to the number of citations received. Node color reflects the relative proportion of incoming vs outgoing influences. Greyed author names have published in HDSR.}
    \label{fig:hdsr_flower}
\end{figure}

% \begin{abstract}
% The abstract should be no more than 250 words.
% \end{abstract}


\vspace*{0.15in}
\hspace{10pt}
  \small	
  \textbf{\textit{Keywords: }} {science of science, information visualization}\\
  \textbf{\textit{Possible HDSR Columns: }} {Meta Data Science, Diving into Data, Bits and Bytes}
% \copyrightnotice
\end{center}
\vspace*{0.15in}

% \section*{Media Summary}
% The Media Summary should be written in plain language to highlight the key messages of
% the article, in ways that can be understood by the general public and cited by the media 
% directly and accurately.  It therefore should avoid technical terms or language designed for
% academic communications. It should not exceed 400 words, and more succinct, the better.



%\section{Proposed Article}
Citations, references, mentions, and hyperlinks leave quantitative traces of how ideas travel through scholarship and knowledge creation. Most assessments of academic impact compress these traces into scalar summaries and rankings. Influence Flower takes a different approach: it is a data-driven visual representation of the incoming and outgoing intellectual influence surrounding a paper, researcher, venue, or organization. The proposed article will use Influence Flower to examine two complementary questions. First, what can a relational visualization reveal that a ranking cannot? Second, what does it take to keep such a scholarly instrument useful when the data, software ecosystem, and operating conditions beneath it continue to change?

\section{Flower as a Visual Metaphor}
The article will begin with the visual metaphor itself, using the HDSR flower in \cref{fig:hdsr_flower} as an example. The flower is built from quantitative scholarly data, but it deliberately does not collapse those data into a single score. Direction distinguishes intellectual influences flowing into and out of the focal entity; arc thickness conveys relative strength; and color helps reveal asymmetry between the two directions. A ranking asks who has more impact. A flower instead invites questions such as where influence comes from, where it travels next, and whether those communities overlap. We will discuss the design choices behind this representation, including what information is gathered, highlighted, and suppressed, and how those choices shape its explanatory value.

This perspective supports qualitative inquiry beyond leaderboards. The 2019 Influence Flower work (\cite{shin2019influence}) demonstrated use cases such as tracing how a researcher's intellectual influences change over a career, identifying papers whose recognition accumulates slowly, summarising the interdisciplinary profile of an institution, and examining how the topical identity of a venue changes over time. These examples motivate a broader argument: scholarly influence is relational, directional, and temporal, and some of those properties are easier to inspect visually than through a scalar ranking. We will position the flower as an exploratory instrument for the science of science rather than as a replacement for quantitative measures of impact.



\section{A Decade of Engineering Evolution}
The second part of the article will tell the system history as a sequence of adaptation pressures rather than as a catalogue of technologies. The conceptual lineage began with static influence visualizations in 2016, while the software project took shape in 2017. Across the following decade, the implementation passed through five distinct eras: rapid prototyping with SQLite and Microsoft Academic Graph (MAG) (\cite{sinha2015overview}); an Elasticsearch-backed system that scaled search and scoring; the 2021 rewrite to a custom memory-mapped citation-graph scoring engine for low-latency computation; a migration from the retired MAG data source to OpenAlex(\cite{priem2022openalex}); and, in 2026, a refactor that removed runtime dependence on external scholarly APIs. For each era, we will use the same narrative structure: what changed, what became too costly or fragile, what adaptation followed, and which parts of the system could remain stable.

The custom graph-scoring engine marks a central transition in this history. Integration with arXiv Labs required influence calculations to move beyond Elasticsearch queries and support arXiv-scale access. The engine stores the citation graph as memory-mapped binary arrays in compressed sparse row form and traverses them using Numba-compiled kernels, allowing fresh flowers to be computed on demand rather than precomputed. The arXiv Labs integration, merged in 2023, also changed the status of the project: a research artifact became a service embedded in other researchers' everyday browsing. Measurements from 2024-2025 show that roughly three quarters of visitors arrived through arXiv referrals, while the integration machinery generated on the order of 750,000 requests per month. Real users turned continued maintenance from an optional virtue into an operational obligation.

Two later transitions make the sustainability problem especially concrete. Microsoft retired the MAG APIs, forcing the system to migrate to OpenAlex. Rather than rewriting the graph engine, the migration introduced an adapter that converted OpenAlex snapshots into the representation expected by the existing graph-scoring pipeline. Then, in 2026, OpenAlex moved to usage-based API pricing. At the traffic generated by the arXiv integration, a metered runtime dependency was no longer practical. The resulting refactor brought names, titles, search, and graph computation under local control using OpenAlex snapshots, the local graph-scoring engine, and self-hosted OpenSearch. The current system has zero runtime dependency on external scholarly APIs and operates over approximately 500 million works, 95 million authors, and 3.04 billion citation edges. The emphasis will be on the architectural decisions, adaptation pressures, and trade-offs that generalize beyond the project.


\section{From Frozen Artifacts to Living Scholarly Infrastructure}

The final part will step back from chronology and ask what this history teaches about sustainable scholarly software. Research papers are designed to become fixed records, and reproducibility often encourages us to freeze code and data alongside them. Public-facing scholarly infrastructure faces the opposite condition: its data sources evolve, APIs disappear or become costly, software dependencies change, traffic grows, and users continue to expect current results. We will contrast the frozen snapshot with living infrastructure and pose the central question of the article: how can living infrastructure be engineered so that change remains inexpensive? Our working answer is that sustainability is not the ability to avoid change, but the ability to absorb inevitable change at a tolerable technical and organizational cost.

We will organize the resulting lessons around four principles. First, design for replaceability: components should be allowed to change without forcing the whole system to change with them. Second, treat external APIs as dependencies with unknown expiration dates, whether the eventual shock is technical, organizational, or economic. Third, optimize for maintenance economics rather than architectural fashion. The memory-mapped graph engine illustrates that memory mapping and the operating system page cache can be a practical data layer when a small academic team values simplicity and predictable operating cost as much as peak throughput. Fourth, make systems cheap to hibernate: research software should be able to survive periods of low maintenance without accumulating enough cost or operational fragility to make revival impossible.

The article will close by returning to the life-cycle metaphor already present in the project. A flower may be visually ephemeral, but the value of a long-lived research system also lies in the "fruits" produced along the way. For Influence Flower these include reusable infrastructure and engineering knowledge; a community of users reached through arXiv Labs; and a way of thinking about scholarly impact as a relational pattern rather than only a leaderboard. The broader lesson is that the flower remained recognizable because its implementation did not remain fixed. Nearly a decade of continuity came from preserving the purpose of the system while allowing almost everything underneath it to be replaced.

\section{Fit for HDSR}

The proposed article sits at the intersection of science of science, information visualization, data engineering, and research-software sustainability. It is not intended as a conventional systems paper or a retrospective release note. Instead, Influence Flower will serve as a concrete case study through which to discuss the scientific value of qualitative visualization and the engineering conditions that allow a public data-science artifact to remain useful over time. We expect the article to be accessible to technical readers and to scholars interested in how computational tools shape the study of science.

% All references should be stored in the file "references.bib"
% Please do not modify anything below this line.
\printbibliography


\end{document}
